% 自 notes/docs/1151/seminar/W09_1109_預習.md（87808f3）由 tools/md_to_tex.py 於 2026-11-09 轉出一次；此後本檔為正本。
\chapter{主要用語對照}

\begin{longtable}[]{@{}
  >{\raggedright\arraybackslash}p{(\linewidth - 6\tabcolsep) * \real{0.3700}}
  >{\raggedright\arraybackslash}p{(\linewidth - 6\tabcolsep) * \real{0.2700}}
  >{\raggedright\arraybackslash}p{(\linewidth - 6\tabcolsep) * \real{0.1300}}
  >{\raggedright\arraybackslash}p{(\linewidth - 6\tabcolsep) * \real{0.2300}}@{}}
\caption{三篇的主要用語。篇目欄 AS 為 Alexander 與 Schauer，FP 為 Ferejohn 與 Pasquino。}\tabularnewline
\toprule\noalign{}
\begin{minipage}[b]{\linewidth}\raggedright
原文
\end{minipage} & \begin{minipage}[b]{\linewidth}\raggedright
中譯
\end{minipage} & \begin{minipage}[b]{\linewidth}\raggedright
篇目
\end{minipage} & \begin{minipage}[b]{\linewidth}\raggedright
出處（頁次）
\end{minipage} \\
\midrule\noalign{}
\endfirsthead
\toprule\noalign{}
\begin{minipage}[b]{\linewidth}\raggedright
原文
\end{minipage} & \begin{minipage}[b]{\linewidth}\raggedright
中譯
\end{minipage} & \begin{minipage}[b]{\linewidth}\raggedright
篇目
\end{minipage} & \begin{minipage}[b]{\linewidth}\raggedright
出處（頁次）
\end{minipage} \\
\midrule\noalign{}
\endhead
\bottomrule\noalign{}
\endlastfoot
judicial supremacy & 司法至上 & AS、Kramer & 標題、456；Kramer 621、623、633 \\
interpretive supremacy & 解釋至上 & AS & 457、459 \\
settlement（for settlement\textquotesingle s sake） & 定紛止爭（為定紛而定紛） & AS & 467、469 \\
first-order / second-order settlement & 一階定紛／二階定紛 & AS & 457、467 \\
content-independent reasons & 與內容無關之理由 & AS & 457 註 11、469 \\
preconstitutional & 先於憲法的 & AS & 459--463 \\
grundnorm / rule of recognition & 基本規範／承認規則 & AS & 460 註 17--18 \\
social fact（acceptance as） & 社會事實（作為社會事實的接受） & AS & 465 \\
Cooper rule & Cooper 規則 & AS & 466 \\
institutional anarchy & 制度上的無政府狀態 & AS & 466、473 \\
asymmetry of authority & 權威的不對稱 & AS & 472 註 50、474 註 51 \\
case-decider / law-maker & 案件裁判者／造法者 & AS & 478--479 \\
dispute settler / constitutional interpreter & 爭端解決者／憲法解釋者 & AS & 479 \\
advisory opinion & 諮詢意見 & AS & 481 \\
stare decisis & 遵循先例 & AS & 477、479 \\
judicial (self-)restraint／deference & 司法自制／尊重 & Kramer & 621、624--625 \\
popular constitutionalism & 人民憲政主義 & Kramer & 622（註 3） \\
Republican／Jeffersonian view & 共和派（Jeffersonian）觀點 & Kramer & 622 \\
Federalist & 聯邦派 & Kramer & 623--626 \\
interpretive authority & 解釋權限（詮釋權威） & Kramer & 622、623、632 \\
prudence／expediency & 審慎／權宜 & Kramer & 626、633 \\
countermajoritarian difficulty & 反多數決困難 & Kramer & 631 \\
interpretive theory／methodology & 解釋理論／解釋方法 & Kramer & 621、632 \\
equilibrium point & 均衡點 & Kramer & 634 \\
principal／agent & 本人／代理人 & Kramer & 625--626 \\
\en{Lochner} era & \en{Lochner} 時代 & Kramer & 630、631 \\
Kelsenian type / Kelsenian model & Kelsen 型（模式） & FP & 21；30；33 \\
burden of legitimation & 合法性負擔 & FP & 21 \\
deliberative expectations & 審議的期待 & FP & 22 \\
public reasons & 公共理由 & FP & 23（Rawls 的語彙） \\
internal deliberation / internally deliberative & 內部審議 & FP & 23；35；36 \\
externally (deliberative) & 外部審議 & FP & 35 \\
parliamentary sovereignty regimes & 議會主權體制 & FP & 28 \\
division of power / separation of powers regimes & 權力分立體制 & FP & 29；33 \\
a priori review & 事前審查（先於公布） & FP & 28；30 \\
abstract review & 抽象審查 & FP & 28；30 \\
concrete (a posteriori) review & 具體（事後）審查 & FP & 31 \\
dispersed review / concentrated review & 分散審查／集中審查 & FP & 29；31；32 \\
bloc de constitutionnalité & 憲法整體 & FP & 28 註 9；34 註 19 \\
negative legislation & 消極立法 & FP & 30 註 13 \\
reason free zone & 無須理由的領域 & FP & 26 \\
countermajoritarian & 反多數決 & FP & 25 \\
\end{longtable}
